\chapter{Il paradigma REST}

Nel capitolo precedente abbiamo sviluppato un'applicazione web completa con
Express. La nostra applicazione si occupava di \textbf{tutto} sul server: gestiva
i dati (con Sequelize), la logica di business (rotte, controller, middleware,
sessioni) e la presentazione (template Pug). Ai browser restava soltanto il
compito di visualizzare le pagine HTML.

\section{I limiti delle web app tradizionali}

Questo approccio è noto come \term{web app tradizionale}. È nato quando i
browser avevano capacità limitate e ai contenuti web si accedeva generalmente
solo tramite browser.

Oggi lo scenario è più complesso:

\begin{itemize}
  \item ai dati possono dover accedere \textbf{altri programmi}, per esempio le
        app mobili;
  \item la tendenza è sfruttare le capacità dei browser moderni per offrire
        un'esperienza più reattiva, cioè le \textit{Single Page App}.
\end{itemize}

\begin{figure}[H]
  \centering
  \begin{tikzpicture}[font=\Lato\scriptsize,
      b/.style={rounded corners=3pt, line width=0.8pt, align=center,
                minimum width=2.4cm, minimum height=0.9cm},
      s/.style={b, draw=primary!65!black, fill=primary!10!white},
      c/.style={b, draw=accent!60!black, fill=accent!10!white},
      a/.style={-{Stealth[length=5pt,width=4pt]}, line width=0.85pt,
                draw=neutral!45!black}]

    \node[c] (br) at (0,0.6)  {Browser};
    \node[c] (na) at (0,-0.9) {App native};

    \node[s] (pr) at (4.4,0.6)  {Presentazione};
    \node[s] (lg) at (7.0,0.6)  {Logica};
    \node[s] (dt) at (9.6,0.6)  {Dati};

    \draw[a] (br) -- (pr);
    \draw[a] (na.east) to[out=0, in=-135] (pr.south west);
    \draw[a] (pr) -- (lg);
    \draw[a] (lg) -- (dt);

    \begin{scope}[on background layer]
      \node[rounded corners=5pt, draw=primary!35!white, line width=0.9pt,
            fit=(pr)(lg)(dt), inner sep=9pt,
            label={[font=\Lato\scriptsize, text=primary!60!black]above:
            applicazione web (server)}] {};
    \end{scope}

    \node[font=\Lato\scriptsize\itshape, text=red-gradient-6!75!black,
          align=center] at (4.7,-1.4)
      {un'app nativa dovrebbe scaricare l'HTML\\ e analizzarlo per estrarre i dati};
  \end{tikzpicture}
  \caption{Nella web app tradizionale i dati sono accessibili solo dentro
           l'HTML.}
  \label{fig:traditional-webapp}
\end{figure}

Le web app tradizionali non sono abbastanza flessibili in questi scenari. Un
altro programma che voglia accedere ai dati dovrebbe:

\begin{enumerate}
  \item scaricare le pagine HTML;
  \item analizzarle per estrarre i dati rilevanti.
\end{enumerate}

Non è efficiente (si trasferiscono moltissimi dati inutili) e non è robusto:
qualunque modifica alle pagine HTML può rompere l'estrazione.

\section{La tendenza attuale}

La tendenza è dividere le applicazioni web in due componenti:

\begin{itemize}
  \item il \term{backend}, responsabile dei dati e della logica di business, sul
        server;
  \item il \term{frontend}, responsabile dell'interfaccia utente, sul client.
        Può essere un'app nativa mobile o desktop, oppure una pagina HTML
        \guillemotleft più intelligente\guillemotright\ che costruisce un'interfaccia interattiva con
        JavaScript, cioè una Single Page App.
\end{itemize}

I due componenti comunicano attraverso Internet. Non è un dettaglio marginale:
secondo il rapporto Akamai sullo stato di Internet (2019), l'\textbf{83\%} del
traffico web è costituito da chiamate ad API.

\dfn{API}{
  Le \term{API} (\textit{Application Programming Interface}) sono il modo in cui
  i programmi comunicano fra loro.}

Come gestire la comunicazione su Internet? Definire a mano un protocollo su
TCP/UDP e aprire socket non è conveniente né interoperabile: se volessimo
integrare $n$ API dovremmo imparare $n$ protocolli diversi. Sono esistite
architetture dedicate (CORBA, Java RMI, SOAP), che in sostanza reinventavano
un'alternativa al Web.

La risposta moderna è più semplice: \textbf{usare i protocolli del Web}.

\section{REST}

\dfn{REST}{
  \term{REST} non è uno standard né un protocollo: è uno \textbf{stile
  architetturale}, cioè un insieme di principi e linee guida che definiscono
  come gli standard web debbano essere usati nelle API. È basato su HTTP e sulle
  URI, e fornisce un'interfaccia comune e coerente basata sull'uso
  \guillemotleft corretto\guillemotright\ di HTTP.}

È lo stile architetturale prevalente per le API web, con un tasso di adozione
del 93,4\%.

\subsection{I fondamenti}

\begin{itemize}
  \item Una API REST permette di interagire con \term{risorse} tramite HTTP.
  \item Tutte le risorse sono associate a una \textbf{URI univoca}.
        \guillemotleft Una risorsa è qualunque cosa sia abbastanza importante da essere
        referenziata come cosa a sé\guillemotright.
  \item Le risorse corrispondono tipicamente (ma non necessariamente) a oggetti
        di dominio persistenti.
  \item I \textbf{verbi HTTP} vanno usati per recuperare o manipolare le
        risorse.
  \item Una API REST dovrebbe essere \textbf{stateless}, cioè non usare sessioni
        lato server.
\end{itemize}

\info{Perché la statelessness migliora la scalabilità}{
  È la stessa ragione vista nel primo capitolo, a proposito di HTTP. Se il
  server non conserva stato fra una richiesta e l'altra, due richieste dello
  stesso client possono essere servite da \textbf{macchine diverse}: per reggere
  più traffico basta aggiungere macchine. Con le sessioni in memoria, invece,
  ogni client resta legato alla macchina che lo ha servito la prima volta.}

\subsection{Representational State Transfer}

Il nome per esteso descrive con precisione il meccanismo:

\begin{itemize}
  \item lo \term{stato dell'applicazione} vive sul \textbf{client};
  \item lo \term{stato della risorsa} vive sul \textbf{server};
  \item i due comunicano trasferendo \term{rappresentazioni} appropriate (JSON,
        XML, \ldots).
\end{itemize}

\begin{affianco}[0.5][c]
\begin{lstlisting}[style=http, name=Scambio REST]
GET /todos HTTP/1.1

[{"todo": "A"}, {"todo": "B"}, ...]
\end{lstlisting}
\risultato
\begin{immagine}[Il client tiene lo stato dell'applicazione,\\
                 il server quello delle risorse.]
\begin{tikzpicture}[    att/.style={rounded corners=3pt, draw=primary!65!black, fill=primary!12!white,
      line width=0.7pt, font=\Lato\scriptsize\bfseries, minimum width=1.8cm,
      minimum height=0.5cm, align=center},
    vita/.style={dashed, draw=neutral!55!white, line width=0.7pt},
    req/.style={-{Stealth[length=5pt,width=4pt]}, line width=0.8pt, draw=primary!70!black},
    res/.style={-{Stealth[length=5pt,width=4pt]}, line width=0.8pt, draw=green-gradient-6},
    msg/.style={mcod, above, inner sep=1.5pt}]
  \node[att] (c) at (0,0) {client};
  \node[att] (s) at (4.4,0) {server};
  \node[mnota, anchor=north] at ([yshift=-1.75cm]c.south) {stato dell'applicazione};
  \node[mnota, anchor=north] at ([yshift=-1.75cm]s.south) {stato delle risorse};
  \draw[vita] (c.south) -- (0,-1.8);
  \draw[vita] (s.south) -- (4.4,-1.8);
  \draw[req] (0,-0.75) -- node[msg]{GET /todos} (4.4,-0.75);
  \draw[res] (4.4,-1.4) -- node[msg]{rappresentazione JSON} (0,-1.4);
\end{tikzpicture}
\end{immagine}
\end{affianco}

\subsection{Formati di interscambio}

I formati più usati sono \term{JSON} e \term{XML}.

\begin{affianco}[0.45]
\codicepiccolo
\begin{lstlisting}[style=json, name=todos.json]
{
  "todos": [
    { "todo": "Learn REST",       "done": false },
    { "todo": "Learn JavaScript", "done": true  }
  ]
}
\end{lstlisting}
\risultato
\codicepiccolo
\begin{lstlisting}[style=html, name=todos.xml]
<?xml version="1.0" encoding="UTF-8" ?>
<root>
  <todos>
    <todo><text>Learn REST</text><done>false</done></todo>
    <todo><text>Learn JavaScript</text><done>true</done></todo>
  </todos>
</root>
\end{lstlisting}
\end{affianco}

JSON si è affermato perché è più compatto e, soprattutto, perché in JavaScript
si trasforma in un oggetto con una sola chiamata.

\subsection{Verbi e URI}

\begin{center}
\renewcommand{\arraystretch}{1.45}
\begin{tabular}{@{}l@{\hspace{1.0cm}}>{\raggedright\arraybackslash}p{8.4cm}@{}}
  \textbf{Verbo e URI} & \textbf{Significato} \\
  \code{GET /todos}           & recupera la lista di tutti gli elementi salvati \\
  \code{GET /todos/<ID>}      & recupera solo l'elemento con quell'identificatore \\
  \code{POST /todos}          & salva un nuovo elemento; i dati sono nel corpo
                                della richiesta \\
  \code{PUT /todos/<ID>}      & sostituisce (o crea) l'elemento con
                                quell'identificatore, usando i dati nel corpo \\
  \code{DELETE /todos/<ID>}   & elimina l'elemento con quell'identificatore \\
\end{tabular}
\end{center}

Qui si vede perché nel primo capitolo insistevamo sulla semantica dei verbi:
in REST è precisamente ciò che rende l'API prevedibile.

\subsection{Risorse annidate}

Talvolta esiste una \textbf{gerarchia} fra le risorse: un'azienda ha zero o più
reparti, che hanno zero o più dipendenti. Progettando un'API è possibile
definire risorse annidate, per esempio \code{GET /seller/\{id\}/reviews} per
elencare tutte le recensioni di un dato venditore.

Va però tenuto presente che ogni risorsa dovrebbe avere \textbf{una sola} URI.

\section{Implementare una API REST}

Le API REST sono concettualmente semplici: basta mettersi in ascolto di
richieste HTTP, gestirle e inviare una risposta. È esattamente quello che
facevamo con la web app tradizionale, con tre differenze:

\begin{itemize}
  \item invece di rispondere con una rappresentazione HTML da disegnare in un
        browser, si usa una rappresentazione \textbf{adatta alle macchine},
        tipicamente JSON;
  \item invece di ricevere input dall'invio di form, li si riceve in una
        rappresentazione adatta alle macchine;
  \item non ci si deve appoggiare alle \textbf{sessioni}: servono altri schemi
        di autenticazione.
\end{itemize}

\section{Autenticazione con token}

Le API REST permettono agli utenti di manipolare risorse, e nella maggior parte
dei casi non vogliamo che chiunque possa farlo.

\begin{itemize}
  \item \term{Autenticazione}: vogliamo che solo gli utenti legittimi accedano
        alle risorse.
  \item \term{Autorizzazione}: vogliamo anche che alcuni utenti accedano solo a
        certe risorse (un dipendente non dovrebbe poter modificare il proprio
        stipendio).
\end{itemize}

Uno schema molto diffuso è basato sui \term{token}:

\begin{enumerate}
  \item il client invia all'API una richiesta con nome utente e password;
  \item l'API valida le credenziali e genera un token;
  \item il token (una stringa) viene restituito al client;
  \item il client deve rimandare il token all'API a \textbf{ogni} richiesta
        successiva che richieda autenticazione, nell'header
        \code{Authorization};
  \item l'API verifica il token prima di rispondere.
\end{enumerate}

\begin{figure}[H]
  \centering
  \begin{tikzpicture}[font=\Lato\small,
      hdr/.style={wtnode, minimum width=2.8cm, minimum height=0.9cm},
      life/.style={draw=neutral!45!white, line width=0.7pt, dashed},
      msg/.style={-{Stealth[length=6pt,width=5pt]}, line width=0.9pt,
                  draw=primary!70!black},
      rsp/.style={-{Stealth[length=6pt,width=5pt]}, line width=0.9pt,
                  draw=accent!70!black},
      bad/.style={-{Stealth[length=6pt,width=5pt]}, line width=0.9pt,
                  draw=red-gradient-6},
      ml/.style={font=\FiraCodeNL\scriptsize, midway, above, inner sep=2.5pt,
                 fill=white, text=neutral!15!black}]

    \def\xb{0} \def\xs{9.0}
    \node[hdr] (B) at (\xb,0) {\textbf{Client}};
    \node[hdr] (S) at (\xs,0) {\textbf{API}};
    \draw[life] (B.south) -- (\xb,-4.4);
    \draw[life] (S.south) -- (\xs,-4.4);

    \draw[bad] (\xb,-0.8) -- (\xs,-0.8) node[ml]{GET /todos (nessun token)};
    \draw[bad] (\xs,-1.5) -- (\xb,-1.5) node[ml]{401 Unauthorized};

    \draw[msg] (\xb,-2.3) -- (\xs,-2.3) node[ml]{POST /auth \{usr, pwd\}};
    \draw[rsp] (\xs,-3.0) -- (\xb,-3.0) node[ml]{200 OK + token};

    \draw[msg] (\xb,-3.7) -- (\xs,-3.7) node[ml]{GET /todos Authorization: Bearer ...};
    \draw[rsp] (\xs,-4.3) -- (\xb,-4.3) node[ml]{200 OK + dati};
  \end{tikzpicture}
  \caption{Lo schema di autenticazione basato su token.}
  \label{fig:token-auth}
\end{figure}

\subsection{JSON Web Token}

\dfn{JWT}{
  \term{JWT} (\textit{JSON Web Token}) è uno standard aperto molto diffuso
  (RFC 7519) che permette di condividere in modo sicuro delle
  \term{asserzioni} (\textit{claim}) fra due parti. Un token JWT è una stringa
  composta da \textbf{tre parti} separate da un punto:
  \code{Header.Payload.Signature}.}

\begin{affianco}[0.6][c]
\codicepiccolo
\begin{lstlisting}[style=plain, name=Un token JWT]
eyJhbGciOiJIUzI1NiIsInR5cCI6IkpXVCJ9
.eyJuYW1lIjoibHVpZ2kiLCJyb2xlIjoiYWRtaW4iLCJleHAiOjE2NzAzOTg0MzJ9
.fopBYrax8wcB7rnPjCcOMc62IT2lJdvyOdyixMWMZAQ
\end{lstlisting}
\risultato
\begin{immagine}[Le tre parti, decodificate: solo la firma\\
                 richiede la chiave segreta.]
\begin{tikzpicture}[
    p/.style={rounded corners=2pt, line width=0.7pt, font=\FiraCodeNL\scriptsize,
      inner sep=3pt, minimum width=1.4cm, anchor=west},
    d/.style={mcod, anchor=west, align=left}]
  \node[p, draw=red-gradient-6, fill=red-gradient-1!40!white] (h) at (0,0) {Header};
  \node[d] at (1.7,0) {\{"alg": "HS256", "typ": "JWT"\}};
  \node[p, draw=accent!60!black, fill=accent!14!white] (y) at (0,-0.75) {Payload};
  \node[d] at (1.7,-0.75) {\{"name": "luigi", "role": "admin",\\ \ "exp": 1670398432\}};
  \node[p, draw=primary!60!black, fill=primary!12!white] (s) at (0,-1.6) {Signature};
  \node[d] at (1.7,-1.6) {HMACSHA256(\ldots, secret\_key)};
\end{tikzpicture}
\end{immagine}
\end{affianco}

\textbf{Header.} Contiene informazioni sul tipo di token e sulla funzione di
hash usata per firmarlo, in formato JSON, codificate in Base64Url.

\begin{lstlisting}[style=json]
{ "alg": "HS256", "typ": "JWT" }
\end{lstlisting}

\textbf{Payload.} È un oggetto JSON contenente le asserzioni. Alcune sono
\textit{registrate} e raccomandate (per esempio \code{exp}, che indica la data
di scadenza del token); altre sono liberamente personalizzabili. Anch'esso è
codificato in Base64Url.

\begin{lstlisting}[style=json]
{ "name": "luigi", "role": "admin", "exp": 1670398432 }
\end{lstlisting}

\error{Base64Url è una codifica, non una cifratura}{
  Header e payload di un JWT si decodificano banalmente: chiunque intercetti il
  token può \textbf{leggerne il contenuto}. Nel payload non vanno quindi mai
  messi dati riservati. Ciò che la firma garantisce è l'\textbf{integrità} (il
  token non è stato alterato), non la riservatezza.}

\textbf{Signature.} Per garantire che i token non siano manomessi o falsificati
si include una firma, ottenuta applicando una funzione di hash crittografica
alla concatenazione di:

\begin{itemize}
  \item l'header del token;
  \item il payload del token;
  \item una \textbf{chiave segreta} nota solo al server che emette il token.
\end{itemize}

\begin{affianco}[0.62][c]
\codicepiccolo
\begin{lstlisting}[style=js]
HMACSHA256(
  base64UrlEncode(header) + "." + base64UrlEncode(payload),
  secret_key
)
\end{lstlisting}
\risultato
\begin{immagine}[Senza la chiave non si può ricalcolare la firma\\
                 di un payload modificato.]
\begin{tikzpicture}[
    b/.style={rounded corners=2pt, line width=0.7pt, font=\FiraCodeNL\scriptsize,
      inner sep=3pt},
    a/.style={-{Stealth[length=4pt,width=3.5pt]}, line width=0.7pt, draw=neutral!45!black}]
  \node[b, draw=red-gradient-6, fill=red-gradient-1!40!white] (h) at (0,0.6) {header};
  \node[b, draw=accent!60!black, fill=accent!14!white] (p) at (0,0) {payload};
  \node[b, draw=neutral!50!black, fill=neutral!10!white] (k) at (0,-0.6) {secret\_key};
  \node[b, draw=neutral!60!black, fill=white, font=\Lato\scriptsize\bfseries] (f) at (2.0,0) {HMAC};
  \node[b, draw=primary!60!black, fill=primary!12!white] (s) at (3.8,0) {signature};
  \foreach \n in {h,p,k} \draw[a] (\n.east) -- (f.west);
  \draw[a] (f) -- (s);
\end{tikzpicture}
\end{immagine}
\end{affianco}

\subsection{Funzioni di hash crittografiche}

\dfn{Funzione di hash crittografica}{
  Una funzione che associa a una stringa binaria di lunghezza arbitraria
  (l'input) una stringa binaria di \textbf{lunghezza fissa} (il \textit{digest}
  o \textit{hash}).}

Le funzioni di hash crittografiche sicure hanno due proprietà interessanti:

\begin{itemize}
  \item sono \textbf{praticamente prive di collisioni}: è di fatto impossibile
        trovare due input diversi che producano lo stesso digest;
  \item sono \textbf{facili da calcolare ma non invertibili}: dato un input è
        immediato calcolarne il digest, ma dato un digest è impraticabile
        risalire all'input che lo ha generato.
\end{itemize}

\begin{center}
\renewcommand{\arraystretch}{1.4}
\begin{tabular}{@{}>{\raggedright\arraybackslash}p{4.0cm}@{\hspace{0.8cm}}>{\raggedright\arraybackslash}p{9.8cm}@{}}
  \textbf{Input} & \textbf{Digest SHA2-256} \\
  \code{Web Technologies}
    & {\FiraCodeNL\scriptsize 72bdf1226df1010dfb47cffaed85f5b35ae15ec9b482388691fa7a662a114113} \\
  \code{Web Technologies!}
    & {\FiraCodeNL\scriptsize 1a6791a2eab52c7c71d3e1225632d866b48a3abad91df15b2d44ab6004877993} \\
\end{tabular}
\end{center}

Un solo carattere di differenza produce un digest completamente diverso: è
questa proprietà a rendere rilevabile qualunque manomissione.

\section{Una API REST con Express}

Riscriviamo la nostra applicazione come API REST, riusando gran parte del codice
Express.

\subsection{Configurazione}

\begin{affianco}[0.55][c]
\begin{lstlisting}[style=js, name=index.js]
const app = express();
const PORT = 3000;

app.use(morgan('dev'));    // log delle richieste
app.use(cors());           // accessibile da ogni origine
app.use(express.json());   // analizza i corpi JSON
\end{lstlisting}
\risultato
\begin{immagine}[Niente sessioni né template: l'API parla JSON.]
\begin{tikzpicture}[
    m/.style={rounded corners=2pt, draw=primary!55!black, fill=primary!8!white,
      line width=0.6pt, font=\FiraCodeNL\scriptsize, minimum width=2.6cm,
      minimum height=0.42cm},
    a/.style={-{Stealth[length=3.5pt,width=3pt]}, line width=0.6pt, draw=neutral!45!black}]
  \foreach \t [count=\i] in {morgan, cors, express.json} \node[m] (m\i) at (0,-\i*0.62) {\t};
  \node[m, fill=accent!12!white, draw=accent!55!black] (r) at (0,-2.5) {router};
  \draw[a] (m1) -- (m2); \draw[a] (m2) -- (m3); \draw[a] (m3) -- (r);
\end{tikzpicture}
\end{immagine}
\end{affianco}

Si noti che sparisce \code{express.urlencoded} (non ci sono più form) e compare
\code{express.json}. Su \code{cors()} torneremo nel capitolo sulla sicurezza.

\subsection{Gestione degli errori}

\begin{affianco}[0.55][c]
\begin{lstlisting}[style=js, name=index.js]
app.use((err, req, res, next) => {
  console.log(err.stack);
  res.status(err.status || 500).json({
    code: err.status || 500,
    description: err.message || "Si e' verificato un errore"
  });
});
\end{lstlisting}
\risultato
\begin{lstlisting}[style=http, name=Risposta d'errore]
HTTP/1.1 404 Not Found
Content-Type: application/json

{"code": 404,
 "description": "Todo non trovato"}
\end{lstlisting}
\end{affianco}

La differenza rispetto alla web app tradizionale è tutta nell'ultima riga:
\code{res.json()} al posto di \code{res.render()}. Anche gli errori devono
essere leggibili da un programma.

\subsection{Autenticazione}

\begin{affianco}[0.6][c]
\codicepiccolo
\begin{lstlisting}[style=js, name=routes/authenticationRouter.js]
authenticationRouter.post("/auth", async (req, res) => {
  let isAuthenticated = await AuthController.checkCredentials(req, res);

  if (isAuthenticated) {
    res.json(AuthController.issueToken(req.body.usr));
  } else {
    res.status(401);
    res.json({ error: "Credenziali non valide. Riprova." });
  }
});

authenticationRouter.post("/signup", (req, res, next) => {
  AuthController.saveUser(req, res).then((user) => {
    res.json(user);
  }).catch((err) => {
    next({ status: 500, message: "Impossibile salvare l'utente" });
  })
});
\end{lstlisting}
\risultato
\begin{immagine}[Il token ottenuto al login accompagna\\ ogni richiesta successiva.]
\begin{tikzpicture}[    att/.style={rounded corners=3pt, draw=primary!65!black, fill=primary!12!white,
      line width=0.7pt, font=\Lato\scriptsize\bfseries, minimum width=1.8cm,
      minimum height=0.5cm, align=center},
    vita/.style={dashed, draw=neutral!55!white, line width=0.7pt},
    req/.style={-{Stealth[length=5pt,width=4pt]}, line width=0.8pt, draw=primary!70!black},
    res/.style={-{Stealth[length=5pt,width=4pt]}, line width=0.8pt, draw=green-gradient-6},
    msg/.style={mcod, above, inner sep=1.5pt}]
  \node[att] (c) at (0,0) {client};
  \node[att] (s) at (4.0,0) {API};
  \draw[vita] (c.south) -- (0,-3.0);
  \draw[vita] (s.south) -- (4.0,-3.0);
  \draw[req] (0,-0.75) -- node[msg]{POST /auth \{usr, pwd\}} (4.0,-0.75);
  \draw[res] (4.0,-1.4) -- node[msg]{200 "eyJhbGci\ldots"} (0,-1.4);
  \draw[req] (0,-2.1) -- node[msg]{GET /todos} node[mcod, below, font=\FiraCodeNL\tiny]{Authorization: Bearer eyJ\ldots} (4.0,-2.1);
  \draw[res] (4.0,-2.85) -- node[msg]{200 [ \ldots ]} (0,-2.85);
\end{tikzpicture}
\end{immagine}
\end{affianco}

\begin{affianco}[0.6][c]
\begin{lstlisting}[style=js, name=controllers/AuthController.js]
import Jwt from "jsonwebtoken";

static issueToken(username) {
  return Jwt.sign({ user: username }, process.env.TOKEN_SECRET, {
    expiresIn: `${24 * 60 * 60}s`
  });
}

static isTokenValid(token, callback) {
  Jwt.verify(token, process.env.TOKEN_SECRET, callback);
}
\end{lstlisting}
\risultato
\begin{immagine}[\code{sign} produce il token, \code{verify} lo controlla.]
\begin{tikzpicture}[
    b/.style={rounded corners=2pt, line width=0.7pt, font=\FiraCodeNL\scriptsize, inner sep=3pt},
    a/.style={-{Stealth[length=4pt,width=3.5pt]}, line width=0.7pt, draw=neutral!45!black}]
  \node[b, draw=accent!60!black, fill=accent!12!white] (p) at (0,0) {\{ user: "luigi" \}};
  \node[b, draw=primary!60!black, fill=primary!10!white] (t) at (3.5,0) {token};
  \draw[a] (p) -- node[mcod, above]{sign} (t);
  \node[b, draw=green-gradient-6, fill=green-gradient-1!45!white] (ok) at (3.5,-1.3) {decodedToken};
  \draw[a] (t) -- node[mcod, right]{verify} (ok);
  \node[mnota, anchor=north] at (ok.south) {oppure \code{err} se scaduto o alterato};
\end{tikzpicture}
\end{immagine}
\end{affianco}

La chiave segreta viene letta da \code{process.env}, cioè dal file \code{.env}
non versionato di cui abbiamo parlato nel capitolo precedente.

\subsection{Middleware di autenticazione}

\begin{affianco}[0.58][c]
\codicepiccolo
\begin{lstlisting}[style=js, name=middleware/auth.js]
export function enforceAuthentication(req, res, next) {
  const authHeader = req.headers['authorization'];
  const token = authHeader?.split(' ')[1];   // "Bearer <token>"

  if (!token) {
    next({ status: 401, message: "Unauthorized" });
    return;
  }

  AuthController.isTokenValid(token, (err, decodedToken) => {
    if (err) {
      next({ status: 401, message: "Unauthorized" });
    } else {
      req.username = decodedToken.user;
      next();
    }
  });
}
\end{lstlisting}
\risultato
\begin{immagine}
\begin{tikzpicture}[    dec/.style={diamond, aspect=2.4, draw=primary!60!black, fill=primary!8!white,
      line width=0.7pt, font=\Lato\scriptsize, inner sep=1pt, align=center},
    esito/.style={rounded corners=3pt, draw=neutral!50!black, fill=white, line width=0.7pt,
      font=\FiraCodeNL\scriptsize, inner sep=3pt, align=center},
    a/.style={-{Stealth[length=4pt,width=3.5pt]}, line width=0.7pt, draw=neutral!45!black}]
  \node[dec] (q1) at (0,0) {c'è il\\ token?};
  \node[dec] (q2) at (0,-1.7) {è valido?};
  \node[esito, draw=merr, text=merr] (e1) at (2.4,0) {401};
  \node[esito, draw=merr, text=merr] (e2) at (2.4,-1.7) {401};
  \node[esito, draw=mok] (ok) at (0,-3.2) {req.username = \ldots\\ next()};
  \draw[a] (q1) -- node[mnota, above]{no} (e1);
  \draw[a] (q2) -- node[mnota, above]{no} (e2);
  \draw[a] (q1) -- node[mnota, right]{sì} (q2);
  \draw[a] (q2) -- node[mnota, right]{sì} (ok);
\end{tikzpicture}
\end{immagine}
\end{affianco}

\subsection{Autorizzazione}

Nella nostra API vogliamo che gli utenti possano vedere e modificare
\textbf{soltanto i propri} elementi. Che cosa succede se un utente invia una
richiesta \code{DELETE} per un elemento che non gli appartiene?

Non basta verificare \textit{chi} è l'utente: bisogna verificare anche che abbia
il \textbf{permesso} di interagire con quella risorsa.

\begin{affianco}[0.62][c]
\codicepiccolo
\begin{lstlisting}[style=js, name=middleware/authz.js]
export async function ensureUsersModifyOnlyOwnTodos(req, res, next) {
  const user   = req.username;
  const todoId = req.params.id;

  const userHasPermission = await AuthController.canUserModifyTodo(user, todoId);

  if (userHasPermission) {
    next();
  } else {
    next({
      status: 403,
      message: "Non hai i permessi per vedere o modificare questa risorsa"
    });
  }
}
\end{lstlisting}
\risultato
\begin{immagine}
\begin{tikzpicture}[    dec/.style={diamond, aspect=2.4, draw=primary!60!black, fill=primary!8!white,
      line width=0.7pt, font=\Lato\scriptsize, inner sep=1pt, align=center},
    esito/.style={rounded corners=3pt, draw=neutral!50!black, fill=white, line width=0.7pt,
      font=\FiraCodeNL\scriptsize, inner sep=3pt, align=center},
    a/.style={-{Stealth[length=4pt,width=3.5pt]}, line width=0.7pt, draw=neutral!45!black}]
  \node[dec] (q) at (0,0) {il todo è\\ dell'utente?};
  \node[esito, draw=mok] (ok) at (0,-1.6) {next()};
  \node[esito, draw=merr, text=merr] (no) at (2.4,0) {403};
  \draw[a] (q) -- node[mnota, right]{sì} (ok);
  \draw[a] (q) -- node[mnota, above]{no} (no);
\end{tikzpicture}
\end{immagine}
\end{affianco}

\begin{affianco}[0.62][c]
\codicepiccolo
\begin{lstlisting}[style=js, name=routes/todoRouter.js]
todoRouter.get("/todos/:id", ensureUsersModifyOnlyOwnTodos, (req, res, next) => {
  TodoController.findById(req).then((item) => {
    if (item)
      res.json(item);
    else
      next({ status: 404, message: "Todo non trovato" });
  }).catch(err => {
    next(err);
  })
});
\end{lstlisting}
\risultato
\begin{immagine}[Una richiesta attraversa tre controlli\\ prima di arrivare ai dati.]
\begin{tikzpicture}[
    m/.style={rounded corners=2pt, draw=primary!55!black, fill=primary!8!white,
      line width=0.6pt, font=\FiraCodeNL\scriptsize, minimum width=3.5cm,
      minimum height=0.42cm},
    a/.style={-{Stealth[length=3.5pt,width=3pt]}, line width=0.6pt, draw=neutral!45!black},
    errx/.style={mcod, anchor=west, text=merr}]
  \node[mcod] (r) at (0,0) {GET /todos/7};
  \node[m] (m1) at (0,-0.7) {enforceAuthentication};
  \node[m] (m2) at (0,-1.4) {ensureUsersModify\ldots};
  \node[m, fill=accent!12!white, draw=accent!55!black] (m3) at (0,-2.1) {findById};
  \node[mcod, text=mok] (ok) at (0,-2.8) {200 \{ \ldots \}};
  \draw[a] (r) -- (m1); \draw[a] (m1) -- (m2); \draw[a] (m2) -- (m3); \draw[a] (m3) -- (ok);
  \node[errx] at ([xshift=4pt]m1.east) {401};
  \node[errx] at ([xshift=4pt]m2.east) {403};
  \node[errx] at ([xshift=4pt]m3.east) {404};
\end{tikzpicture}
\end{immagine}
\end{affianco}

\info{401 contro 403}{
  Sono due cose diverse e vanno usate con precisione. \code{401 Unauthorized}
  significa \guillemotleft non so chi sei\guillemotright: manca l'autenticazione, o il token non è
  valido. \code{403 Forbidden} significa \guillemotleft so chi sei, ma non puoi\guillemotright:
  l'autenticazione c'è, mancano i permessi.}

\section{OpenAPI}

\dfn{OpenAPI}{
  \term{OpenAPI} (già noto come Swagger) è uno standard aperto e formale per
  \textbf{descrivere} le API HTTP.}

Perché serve?

\begin{itemize}
  \item \textbf{Standardizzazione}: garantisce pratiche di documentazione
        coerenti.
  \item \textbf{Documentazione}: si può usare per generare automaticamente
        documentazione completa.
  \item \textbf{Generazione di codice}: permette di generare automaticamente
        librerie client e scheletri di server.
  \item \textbf{Leggibile da macchine e da persone}: abilita la scoperta
        automatica e altro ancora.
\end{itemize}

È lo standard industriale più adottato.

\subsection{La specifica}

Le descrizioni OpenAPI sono documenti di testo strutturato: ogni documento
rappresenta un oggetto JSON, in formato JSON o YAML. Descrivono la versione
della specifica (\code{openapi}), l'API (\code{info}) e gli endpoint
(\code{paths}).

\begin{lstlisting}[style=plain, name=openapi.yaml]
openapi: 3.1.0
info:
  title: A minimal specification
  version: 0.0.1
paths: {}   # nessun endpoint definito
\end{lstlisting}

Gli endpoint si chiamano \term{path}. Per ciascuno si specificano i metodi
supportati e, per ciascun metodo, i parametri di ingresso attesi, l'eventuale
corpo della richiesta e le possibili risposte.

\begin{affianco}[0.5][c]
\begin{lstlisting}[style=plain, name=openapi.yaml]
/auth:
  post:
    description: Authenticate user
    requestBody:
      description: user to authenticate
      required: true
      content:
        application/json:
          schema:
            type: object
            properties:
              usr:
                type: string
                example: Kyle
              pwd:
                type: string
                example: p4ssw0rd
    responses:
      200:
        description: User authenticated
      401:
        description: Invalid credentials
\end{lstlisting}
\risultato
\begin{immagine}[Come la stessa specifica appare in Swagger UI.]
\begin{tikzpicture}[
    bar/.style={rounded corners=3pt, draw=green-gradient-6, fill=green-gradient-1!35!white,
      line width=0.7pt, minimum width=6.0cm, minimum height=0.6cm, anchor=north west},
    verb/.style={rounded corners=2pt, fill=green-gradient-6, text=white,
      font=\Lato\scriptsize\bfseries, inner sep=3pt},
    r/.style={font=\FiraCodeNL\scriptsize, anchor=west}]
  \node[bar] (b) at (0,0) {};
  \node[verb, anchor=west] (v) at ([xshift=4pt]b.west) {POST};
  \node[r] at ([xshift=4pt]v.east) {/auth};
  \node[mnota, anchor=east] at ([xshift=-4pt]b.east) {Authenticate user};
  \node[mnota, anchor=north west, font=\Lato\scriptsize\bfseries] at (0.1,-0.75) {Request body \textcolor{merr}{* required}};
  \node[r, anchor=north west] at (0.1,-1.15) {\{ "usr": "Kyle", "pwd": "p4ssw0rd" \}};
  \node[mnota, anchor=north west, font=\Lato\scriptsize\bfseries] at (0.1,-1.65) {Responses};
  \node[r, anchor=north west] at (0.1,-2.05) {200 \quad User authenticated};
  \node[r, anchor=north west] at (0.1,-2.45) {401 \quad Invalid credentials};
\end{tikzpicture}
\end{immagine}
\end{affianco}

\subsection{A che cosa serve una specifica}

\begin{center}
\renewcommand{\arraystretch}{1.4}
\begin{tabular}{@{}l@{\hspace{1.2cm}}>{\raggedright\arraybackslash}p{8.6cm}@{}}
  \textbf{Uso} & \textbf{Che cosa se ne ottiene} \\
  Supporto alla progettazione & rilevare errori e incoerenze prima di scrivere
                                codice \\
  Generazione di codice       & scheletri di server, SDK per i client, test \\
  Documentazione              & documentazione interattiva generata
                                automaticamente \\
\end{tabular}
\end{center}

OpenAPI è supportato da diversi strumenti open source maturi: \textbf{Swagger
Editor}, \textbf{Swagger Codegen} e \textbf{Swagger UI}, provabili anche
direttamente nel browser su \urlx{https://editor.swagger.io/}.

La generazione di codice lato server può ovviamente produrre solo scheletri, ma
supporta ASP.NET, Go, Java, JaxRS, Micronaut, Kotlin, Node.js, Python (Flask),
Scala e Spring. Lato client permette di generare SDK in decine di linguaggi
diversi.

\info{Sviluppo guidato dalla specifica}{
  Spesso si comincia a lavorare a una API \textbf{definendone la specifica}
  prima di scrivere qualunque codice: è il cosiddetto \textit{OpenAPI-driven
  development}. Oltre a sfruttare al massimo la generazione di codice, promuove
  l'\textbf{indipendenza} fra i diversi gruppi coinvolti (front-end, back-end,
  qualità): la definizione dell'API li tiene allineati e permette di lavorare in
  parallelo.}

\subsection{Generare la specifica dal codice}

Nel corso si usano due pacchetti Node:

\begin{itemize}
  \item \code{swagger-jsdoc}, che genera automaticamente le specifiche OpenAPI a
        partire da annotazioni nel codice sorgente;
  \item \code{swagger-ui-express}, che serve da Express la documentazione
        Swagger UI generata automaticamente.
\end{itemize}

\begin{lstlisting}[style=shell]
npm install swagger-jsdoc
npm install swagger-ui-express
\end{lstlisting}

\begin{affianco}[0.58][c]
\codicepiccolo
\begin{lstlisting}[style=js, name=index.js]
const swaggerSpec = swaggerJSDoc({
  definition: {
    openapi: '3.1.0',
    info: {
      title: 'To-do List REST API',
      version: '1.0.0',
    },
  },
  apis: ['./routes/*Router.js'],   // file contenenti le annotazioni
});

// Swagger UI sara' disponibile su /api-docs
app.use('/api-docs', swaggerUI.serve, swaggerUI.setup(swaggerSpec));
\end{lstlisting}
\risultato
\begin{browser}[localhost:3000/api-docs]
  {\Lato\bfseries\Large To-do List REST API}\ {\Lato\scriptsize\colorbox{neutral!30!white}{1.0.0}}\\[6pt]
  \begin{tikzpicture}
    \node[rounded corners=3pt, draw=green-gradient-6, fill=green-gradient-1!35!white,
          line width=0.7pt, minimum width=5.4cm, minimum height=0.55cm] (b) at (0,0) {};
    \node[rounded corners=2pt, fill=green-gradient-6, text=white,
          font=\Lato\scriptsize\bfseries, inner sep=3pt, anchor=west] (v) at ([xshift=4pt]b.west) {POST};
    \node[font=\FiraCodeNL\scriptsize, anchor=west] at ([xshift=4pt]v.east) {/auth};
  \end{tikzpicture}
\end{browser}
\wtnota{La documentazione interattiva generata\\ dalle annotazioni nei router.}
\end{affianco}

\begin{lstlisting}[style=js, name=routes/authenticationRouter.js]
/**
 * @swagger
 * /auth:
 *   post:
 *     description: Authenticate user
 *     requestBody:
 *       content:
 *         application/json:
 *           schema:
 *             type: object
 *             properties:
 *               usr:
 *                 type: string
 *               pwd:
 *                 type: string
 */
authenticationRouter.post("/auth", async (req, res) => { /* ... */ });
\end{lstlisting}

\warning{Che cosa serve sapere all'esame}{
  Per il corso è richiesto sapere \textbf{che cos'è} lo standard OpenAPI,
  \textbf{perché} è stato introdotto e \textbf{a che cosa} può servire. Non è
  richiesto saper scrivere file di specifica OpenAPI da zero.}

\section{Riferimenti}

\begin{itemize}
  \item \textbf{The Little Book on REST Services}, Kenneth Lange. \\
        \urlx{https://www.kennethlange.com/books/The-Little-Book-on-REST-Services.pdf}
  \item \textbf{API design guide} --- Google Cloud. \\
        \urlx{https://cloud.google.com/apis/design} \\
        Parti rilevanti: \textit{Resource-oriented design}, \textit{Resource
        names}.
  \item \textbf{REST API Checklist}, Kenneth Lange. \\
        \urlx{https://kennethlange.com/rest-api-checklist/}
  \item \textbf{OpenAPI} --- sito ufficiale. \\
        \urlx{https://www.openapis.org/} \\
        \urlx{https://learn.openapis.org/} \\
        \urlx{https://spec.openapis.org/oas/v3.1.0}
  \item \textbf{Esempi con altri framework}: \\
        \urlx{https://github.com/luistar/jakartaee-rest-example} (JakartaEE) \\
        \urlx{https://github.com/luistar/rest-service-pizzashop} (Spring) \\
        \urlx{https://github.com/marciovrl/fastapi} (FastAPI, Python) \\
        Non sono richiesti per il corso, ma è istruttivo vedere framework che
        adottano un approccio basato su annotazioni.
\end{itemize}
